% ============================================================
% ch08.tex —— 第 8 章 对称动物园：群表
% ============================================================
\chapter{对称动物园：群表}

上一章我们确认"对称动作构成群"。这一章把这个群养大：给它编花名册、画全家福，然后干一件大事——把一张牌桌上的洗牌，和一张三角形卡片，按在\textbf{同一张}乘法表上。到那时你会真正明白"抽象"这个词值多少钱。

\section{花名册：数对称}

\begin{kaiwei}
等边三角形有 $6$ 种复原动作（上一章数过）。正方形有几种？正五边形呢？正 $n$ 边形呢？先猜一个公式，再看答案。
\end{kaiwei}

正方形的账仔细算：\textbf{旋转类}——$0^\circ, 90^\circ, 180^\circ, 270^\circ$ 四种（$360^\circ$ 算 $0^\circ$）；\textbf{翻转类}——正方形有 $4$ 条对称轴（$2$ 条对角线、$2$ 条对边中线），沿每条翻一次，$4$ 种。合计 $8$ 种。一般规律水落石出：
\[
\text{正 } n \text{ 边形的复原动作} = \underbrace{n \text{ 种旋转}}_{0, \tfrac{360^\circ}{n}, \dots} + \underbrace{n \text{ 种翻转}}_{n \text{ 条对称轴}} = 2n \text{ 种} .
\]
验证：三角形 $2\times3 = 6$ \checkmark，正方形 $8$ \checkmark，正五边形 $10$，正六边形 $12$。你每天路过的人行道方砖（正方形地砖），它的对称群就有 $8$ 个成员。

\begin{faming}
\textbf{土名}：复原动作全家福 $\to$ \textbf{正式名}：正 $n$ 边形的\textbf{对称群}，记号 $D_n$，专业名\textbf{二面体群}（dihedral group），大小 $2n$。它是你能亲手摸到的最丰富的群家族——纸片、地砖、雪花（$D_6$）、海星（$D_5$）都是它的会员。
\end{faming}

\section{全家福：群表（凯莱表）}

给三角形的 $6$ 个动作起代号（这套记号第 9、12 章还要用，值得记牢）：
\begin{itemize}
\item $e$：什么都不做（单位元）；
\item $r$：绕中心转 $120^\circ$（rotate）；
\item $r^2$：转 $240^\circ$（连转两次 $120^\circ$，所以写成 $r$ 的平方）；
\item $s$：沿某条固定对称轴翻转（swap/reflect）；
\item $sr$：先转 $120^\circ$ 再翻转（按"先右后左"的约定，$sr$ = 先做 $r$，再做 $s$）；
\item $sr^2$：先转 $240^\circ$ 再翻转。
\end{itemize}
定义乘法"行 $\circ$ 列 $=$ 先做列动作，再做行动作"，把 $36$ 个乘积全列出来，就得到三角形的\textbf{群表}：

\begin{center}
\renewcommand{\arraystretch}{1.3}
\begin{tabular}{c|cccccc}
$\circ$ & $e$ & $r$ & $r^2$ & $s$ & $sr$ & $sr^2$ \\
\hline
$e$    & $e$    & $r$    & $r^2$  & $s$    & $sr$    & $sr^2$ \\
$r$    & $r$    & $r^2$  & $e$    & $sr^2$ & $s$     & $sr$   \\
$r^2$  & $r^2$  & $e$    & $r$    & $sr$   & $sr^2$  & $s$    \\
$s$    & $s$    & $sr$   & $sr^2$ & $e$    & $r$     & $r^2$  \\
$sr$   & $sr$   & $sr^2$ & $s$    & $r^2$  & $e$     & $r$    \\
$sr^2$ & $sr^2$ & $s$    & $sr$   & $r$    & $r^2$   & $e$    \\
\end{tabular}
\end{center}

这张表叫\Keyword{群表}（正式名：凯莱表，Cayley table——凯莱 1854 年就是靠"把运算写成表"这一步，把群从具体运算中解放出来的）。\textbf{请不要背它}。正确姿势：\textbf{剪一个纸三角形，三个顶点写上 $1,2,3$，随手抽格子验算}。

示范两个格子：

\textbf{格子 $(r \text{ 行},\, s \text{ 列}) = sr^2$}，即"$r\circ s$：先翻转 $s$，再旋转 $r$"。纸片操作：初始顶点 $1$ 在正上方。翻转 $s$（比如沿过顶点 $1$ 的轴）：$1$ 不动，$2\leftrightarrow3$ 互换。再转 $120^\circ$：三个顶点整体转一格。最终格局——对照表格，它等于"先转 $240^\circ$ 再翻转"，即 $sr^2$。\checkmark

\textbf{格子 $(s \text{ 行},\, r \text{ 列}) = sr$}，即"$s\circ r$：先转 $120^\circ$，再翻转"。同样的纸片，顺序换一下，落点\textbf{不同}——这引出本章后半场的大新闻。

验算靠纸片，但\textbf{填表}其实可以纯代数完成，两个小技巧就够：
\begin{itemize}
\item \textbf{转三圈回家}：$r^3 = e$（转三次 $120^\circ$ 回原样）。这解释了 $e$ 行下面的 $r, r^2, e$ 循环；
\item \textbf{翻两次回家，且翻转把转向倒过来}：$s^2 = e$；以及关键关系 $s\,r\,s = r^{-1} = r^2$（先翻、再转、再翻 $=$ 反着转——想象在镜子里看旋转：顺时针变逆时针！）。有了 $r^3 = s^2 = e$ 和 $srs = r^{-1}$ 这三条"宪法"，整张 $6\times6$ 表能徒手推出来。
\end{itemize}

\section{群表的铁律：每行每列都是全排列}

盯着表看三十秒，找规律：\textbf{每一行、每一列，都是 $e, r, r^2, s, sr, sr^2$ 六个元素的某种重新排列——没有任何重复、没有任何遗漏}。这不是三角形的巧合，而是群的结构性铁律：

\begin{dingli}[群表无重复]
群 $(G, *)$ 的群表中，任意一行的 $|G|$ 个格子两两不同（从而恰是全体元素的一个排列）；列同理。
\end{dingli}
\begin{proof}
看第 $a$ 行。若有两格重复，$a*x = a*y$（$x \ne y$），由第 7 章的可约去律（左乘 $a^{-1}$），$x = y$——矛盾。所以一行 $|G|$ 个格子互不相同；而格子只能填 $G$ 的元素（封闭性），有限个互不相同的东西恰好取遍全体。列同理（右约去）。
\end{proof}

回头看一眼你背了多年的\textbf{九九乘法表}：它的"模 $7$ 版"（第 7 章习题里 $1$--$6$ 的乘法表）恰好满足这条铁律——因为 $Z_7$ 去零是个群。而完整的九九表（$1$--$9$）\textbf{不}满足：$2\times3 = 1\times6 = 6$，重复了——因为 $Z_{10}$ 去零不是群（$10$ 不是素数，有人没倒数）。\textbf{一张乘法表合不合格，看它像不像"数独"就知道了}。这给了你一个终身的检查技巧：将来任何场合见到"运算表"，扫一眼行有没有重复，立刻知道它离群有多远。

\section{大新闻：次序敏感（不交换的运算）}

九九表里 $3\times5 = 5\times3$，行行对称。但看三角形群表的两个格子：
\[
(r \text{ 行}, s \text{ 列}) = sr^2, \qquad\text{而}\qquad (s \text{ 行}, r \text{ 列}) = sr .
\]
$sr^2 \ne sr$（表里它们是不同元素）。也就是说
\[
r\circ s \;\ne\; s\circ r :
\qquad \textbf{先翻转再旋转} \ne \textbf{先旋转再翻转}.
\]
纸片立刻验证（上一节你已经各转过一遍，落点不同）。这是你在本书遇到的第一个\textbf{不遵守交换律}的运算。加法、乘法、$\heartsuit$ 全都交换，于是我们以为交换是天经地义——错。日常生活其实早有预警：\textbf{先穿袜子再穿鞋}，和\textbf{先穿鞋再穿袜子}，结果大不相同；\textbf{先腌后炸}与\textbf{先炸后腌}是两道菜。对称动作的世界里，次序是生命线。

\begin{faming}
交换的群叫\textbf{阿贝尔群}（Abelian group，纪念挪威数学家阿贝尔——他证明了五次方程一般没有求根公式，第 12 章出场）；不交换的群就叫\textbf{非交换群}或非阿贝尔群。数学从此分成两个大区：阿贝尔区温顺如初等代数（"移项"随便移），非阿贝尔区暗流涌动（移项要分左右）。三元素以上的对称群全是不交换的——\textbf{你想感受非交换，剪张纸片就够了}。
\end{faming}

\section{动物园的巨型居民}

对称动物园越逛越大，最后两件镇馆之宝必须点名：

\begin{itemize}
\item \textbf{魔方}。每个转法（转某层 $90^\circ$）是一个动作，全部转法的复合构成一个群，大小为
\[
43{,}252{,}003{,}274{,}489{,}856{,}000 \;\approx\; 4.3\times10^{19}.
\]
两个直接推论：其一，\textbf{"打乱后总能复原"是群论的定理}——每个动作都有逆元，原路返回即可；其二，高手背的"公式"就是群表里精心挑选的短路径（把任意打乱状态归约到少数"标准态"，再查表收官）。2010 年，借助计算机穷举（把 $4.3\times10^{19}$ 状态按对称性压缩后搜索），人类证明了"\textbf{上帝之数}是 $20$"：任何打乱的魔方，$20$ 步内必能复原。
\item \textbf{一副扑克}。$54$ 张牌的每种排列是一个"重排"（置换），重排之间可以复合——构成 $54! \approx 2.3\times10^{71}$ 阶的群。这个数什么概念？假设全人类 $80$ 亿人从宇宙诞生起每秒各洗一次牌，到今天总洗牌次数约 $10^{27}$ 次——离 $10^{71}$ 还差四十多个数量级。推论：\textbf{一次充分随机的洗牌，产出的排列几乎必然是这副牌自出厂以来从未出现过的}——你手里握着宇宙级的孤品。
\end{itemize}

更妙的是，这两位巨无霸的\textbf{局部}就是你手里的纸片群：魔方只转一个面的四种转法，复合结构含着 $D_4$ 的影子；扑克三张牌的洗牌（下一节）就是 $D_3$ 本尊。

\section{认亲大会：三角形与三张牌}

拿三张扑克，编号 $1, 2, 3$。"洗牌"就是重新排列。全部排列方式 $3! = 6$ 种。挑两种"基本洗法"：
\begin{itemize}
\item \textbf{轮换} $\rho$：$1\to2, 2\to3, 3\to1$（每张牌往后挪一位，末尾绕回开头）；
\item \textbf{对换} $\sigma$：交换 $1, 2$ 两张，$3$ 不动。
\end{itemize}
任何洗法都能由这两个基本洗法复合出来（试：$\rho^2$ 是反着轮换；$\sigma\rho$ 是"先轮换再交换前两张"……）。复合作为乘法，四项体检全部通过（请自行过一遍，与第 7 章对称群的论证逐字相同）——三张牌的洗法构成一个 $6$ 元群，记 $\{e, \rho, \rho^2, \sigma, \sigma\rho, \sigma\rho^2\}$。

列出它的群表，与三角形对称群表并排放好——\textbf{逐格相同}（把 $\rho$ 换成 $r$、$\sigma$ 换成 $s$，一字不差）。两个看起来毫无血缘的对象——几何的翻转旋转、牌桌的排列——共享同一张乘法表。

\begin{faming}
名字不同、表格相同的两个群，称为\textbf{同构}（isomorphic，希腊词根"同形"，土名：失散多年的双胞胎），记作 $G \cong H$。发现一个同构，是抽象代数里最快乐的事件之一：\textbf{它意味着两个世界的一切群论性质完全共享}——一边证明的定理，另一边免费领取。三角形的对称群 $\cong$ 三张牌的洗牌群（这个群名字叫 $S_3$，三元素置换群），今后你只需要研究\textbf{一个}，就等于研究了\textbf{两个}。
\end{faming}

序章烧脑时刻你已经预演过一次同构（$\heartsuit$ 群 $\cong$ 整数加法群：换元 $u = a+1$ 即可翻译）。"透过现象看结构"——\textbf{同构}就是"结构相同"的正式发音，这就是"抽象"二字的全部含义，也是它全部的报酬。

\begin{dongshou}
\begin{itemize}
  \yisi 在三角形群表里找出每个元素的逆元（"$\circ$ 后等于 $e$"的搭档）：$e^{-1} = e$，$r^{-1} = r^2$，$s^{-1} = s$，请补全 $sr, sr^2$ 的。数一数"自己就是自己逆元"的元素有几个（答案 $4$：$e, s, sr, sr^2$——翻转全是"翻了再翻就回来"的脾气）。
  \yier 剪一个正方形纸片，顶点标 $1,2,3,4$。实验比较"先沿对角线翻转，再转 $90^\circ$"与"先转 $90^\circ$，再沿同一条对角线翻转"，画出两次的顶点落位图。它们相等吗？
  \yier 用"宪法三关系"（$r^4 = e$，$s^2 = e$，$srs = r^{-1}$）推出正方形群 $D_4$（$8$ 元）的完整群表。推 $3\times3 = 9$ 个格子后停下来休息，感受"宪法生全表"的杠杆效应。
  \yier 三张牌的洗牌群 $S_3$ 中，$\sigma\rho$ 与 $\rho\sigma$ 相等吗？（用"先右后左"逐张追踪：$\sigma\rho$ = 先 $\rho$ 后 $\sigma$；$\rho\sigma$ = 先 $\sigma$ 后 $\rho$。落位不同——非交换在牌桌上原样上演。）
\end{itemize}
\end{dongshou}

\begin{shaonao}
\textbf{两元群的世界垄断案}。设一个群只有两个元素 $\{e, a\}$，$a \ne e$。问：$a\circ a$ 只能是什么？——不能是 $a$（否则该行出现两个 $a$，违反"群表无重复"）；不能是别的（没有别的元素了）；所以 $a\circ a = e$。于是两元群的群表\textbf{全世界仅此一种}，任何两元群互相同构。给它点名：电灯开关（开/关，按两下回原样）、硬币（正/反，翻两次回原面）、"乘 $1$ 或乘 $-1$"——全是它的化身。请再给这个"宇宙唯一两元群"找两个新化身。\textbf{同构思维的第一课：小世界里，结构可以被完全数清楚}。
\end{shaonao}

\begin{chengshi}
群表能完全刻画小群，但画不了大群（$S_4$ 的表就 $24\times24 = 576$ 格）——研究大群靠"生成元 + 关系"（本章的"宪法"）和第 9 章的子群切割术，这是大学群论课的主线。"上帝之数 $20$"的证明用了大量计算（约 35 CPU 年），数学界公认其为定理，但"计算机辅助证明"的哲学地位至今偶有争论——顺便见识一下数学的自我审视。
\end{chengshi}

动物园逛完，你已经能认出对称、数出成员、看穿同构。下一章是篇二的核心：拿一个群，\textbf{切蛋糕}。切出来的定理将一次还清两笔欠款——费马小定理的完整证明，以及"素数世界里周期为何整除 $p-1$"的全部真相。
